\documentclass[11pt,a4paper]{article}

% ---------------------------------------------------------
% PREAMBLE (stable, warning-free, Overleaf-safe)
% ---------------------------------------------------------
\usepackage[utf8]{inputenc}
\usepackage[T1]{fontenc}
\usepackage{lmodern}
\usepackage{amsmath, amssymb, amsfonts}
\usepackage{bm}
\usepackage{geometry}
\usepackage{graphicx}
\usepackage{hyperref}
\usepackage{cite}
\usepackage{authblk}
\usepackage{physics}
\usepackage{mathtools}
\usepackage{textgreek}

\geometry{margin=2.4cm}

\title{\textbf{Companion LXXV: The Fisher Kernel for Persistence Diagrams in the Relaxation-Driven Cyclic Cosmology}}
\author{Michael Lehmann \\ \texttt{mi.lehmann@gmx.de}}
\date{April 2026}

\begin{document}
\maketitle

\begin{abstract}
The Fisher kernel provides a principled method for combining information geometry 
with reproducing kernel Hilbert space (RKHS) methods. Applied to persistence 
diagrams, it yields feature maps that incorporate both the geometric structure of 
the underlying statistical manifold and the linear structure of kernel methods. 
Building on the Fisher information geometry introduced in Companion LXXIV, this 
work develops the Fisher kernel for persistence diagrams in the Relaxation-Driven 
Cyclic Cosmology (RDCC). The relaxation mechanism modifies the scale-dependent 
evolution of the gravitational potential, inducing characteristic changes in the 
Fisher score and the resulting kernel. The Fisher kernel provides a powerful tool 
for classification, regression, and model discrimination in weak-lensing topology, 
bridging information geometry and RKHS-based topological data analysis.
\end{abstract}
% ---------------------------------------------------------
\section{Introduction}

Persistence diagrams encode the birth and death of topological features in weak-
lensing convergence fields. Previous Companions developed Hilbert-space embeddings, 
kernel PCA, maximum mean discrepancy (MMD), and Fisher information geometry as 
tools for analyzing and distinguishing the Relaxation-Driven Cyclic Cosmology 
(RDCC) from ΛCDM.

The Fisher kernel provides a natural synthesis of these approaches. It maps each 
persistence diagram to a feature vector given by the Fisher score, which measures 
the sensitivity of the log-likelihood to changes in cosmological parameters. The 
inner product of these score vectors defines a kernel that incorporates both 
information geometry and RKHS structure. This kernel is particularly well-suited 
for classification, regression, and hypothesis testing in cosmological topology.
% ---------------------------------------------------------
\section{Fisher Score and Fisher Kernel}

Given a statistical model $p(D|\theta)$ for persistence diagrams $D$, the Fisher 
score is

\begin{equation}
    U_{\theta}(D)
    = \partial_{\theta} \log p(D|\theta).
\end{equation}

The Fisher information matrix is

\begin{equation}
    I(\theta)
    = \mathbb{E}
      \left[
        U_{\theta}(D) U_{\theta}(D)^{\top}
      \right].
\end{equation}

The Fisher kernel is defined as

\begin{equation}
    k_{F}(D_{1}, D_{2})
    = U_{\theta}(D_{1})^{\top}
      I(\theta)^{-1}
      U_{\theta}(D_{2}).
\end{equation}

This kernel incorporates both the curvature of the statistical manifold and the 
local sensitivity of persistence diagrams to parameter changes.
% ---------------------------------------------------------
\section{RDCC Modification of Fisher Scores}

In RDCC, the convergence evolves as

\begin{equation}
    \kappa_{\mathrm{RDCC}}(\ell)
    = \kappa(\ell)
      \left[
        1 - \chi_{\kappa}
        \left(\frac{\ell}{\ell_{\mathrm{rel}}}\right)^{\alpha}
      \right].
\end{equation}

This modifies the distribution of persistence diagrams:



\[
    p_{\mathrm{RDCC}}(D|\theta)
    = p(D|\theta)
      \left[
        1 - \chi_{p}
        \left(\frac{\ell}{\ell_{\mathrm{rel}}}\right)^{\alpha}
      \right].
\]



The Fisher score becomes

\begin{equation}
    U^{\mathrm{RDCC}}_{\theta}(D)
    = U_{\theta}(D)
      \left[
        1 - \chi_{U}
        \left(\frac{\ell}{\ell_{\mathrm{rel}}}\right)^{\alpha}
      \right].
\end{equation}
% ---------------------------------------------------------
\section{The Fisher Kernel in RDCC}

The RDCC Fisher kernel is

\begin{equation}
    k_{F,\mathrm{RDCC}}(D_{1}, D_{2})
    = k_{F}(D_{1}, D_{2})
      \left[
        1 - \chi_{F}
        \left(\frac{\ell}{\ell_{\mathrm{rel}}}\right)^{\alpha}
      \right].
\end{equation}

RDCC induces:

\begin{itemize}
    \item enhanced kernel similarity for large-scale topological modes,
    \item suppressed similarity for small-scale modes,
    \item a characteristic turnover near $\ell_{\mathrm{rel}}$,
    \item modified Fisher curvature of the topological manifold,
    \item altered alignment of Fisher score vectors across redshift.
\end{itemize}
% ---------------------------------------------------------
\section{Applications to Cosmological Topology}

The Fisher kernel enables:

\begin{itemize}
    \item classification of RDCC vs. ΛCDM topological ensembles,
    \item regression of RDCC parameters from persistence diagrams,
    \item kernel PCA in Fisher space,
    \item Fisher-kernel MMD tests,
    \item information-geometric feature extraction for weak-lensing surveys.
\end{itemize}

These applications provide a powerful bridge between information geometry, kernel 
methods, and cosmological inference.
% ---------------------------------------------------------
\section{Conclusion}

The Fisher kernel provides a principled method for combining information geometry 
with RKHS-based analysis of persistence diagrams. In the Relaxation-Driven Cyclic 
Cosmology, the relaxation mechanism modifies the Fisher score and the resulting 
kernel in a scale-dependent manner, producing distinctive signatures in weak-
lensing topology. The present Companion establishes the foundation for future 
studies of Fisher-kernel methods, parameter inference, and information geometry in 
cosmological topology.

\bibliographystyle{apsrev4-2}
\nocite{*}
\bibliography{references}


\end{document}
